\documentclass[12pt, a4paper]{article}
\usepackage{../../../myconfig} % Gọi file cấu hình từ ngoài gốc vào


% ==========================================
% BẮT ĐẦU NỘI DUNG TÀI LIỆU
% ==========================================
\begin{document}

% Truyền 3 tham số: {Nhãn cam} {Chữ to chính giữa} {Chữ ở góc phải Header}
\titlebox{Chủ đề: Hàm số}{Lượng giác}{Hàm số: Lượng giác}

% --- CÂU 1 ---
\questionLinesManual{0}{1}{Trong các công thức sau, công thức nào sai?}{
    \begin{tasks}(2)
        \task \( \cos 2a = \cos^2 a - \sin^2 a \)
        \task \( \cos 2a = \cos^2 a + \sin^2 a \)
        \task \( \cos 2a = 2\cos^2 a - 1 \)
        \task \( \cos 2a = 1 - 2\sin^2 a \)
    \end{tasks}
}

% --- CÂU 2 ---
\questionLinesManual{0}{2}{Chọn đáp án đúng.}{
    \begin{tasks}(2)
        \task \( \sin 2x = 2\sin x \cos x \)
        \task \( \sin 2x = \sin x \cos x \)
        \task \( \sin 2x = 2\cos x \)
        \task \( \sin 2x = 2\sin x \)
    \end{tasks}
}

% --- CÂU 3 ---
\questionLinesManual{0}{3}{Phương trình \( \sin x = 1 \) có một nghiệm là}{
    \begin{tasks}(4)
        \task \( x = \pi \)
        \task \( x = -\dfrac{\pi}{2} \)
        \task \( x = \dfrac{\pi}{2} \)
        \task \( x = \dfrac{\pi}{3} \)
    \end{tasks}
}

% --- CÂU 4 ---
\questionLinesManual{0}{4}{Nghiệm của phương trình \( \cos x = -\dfrac{1}{2} \) là}{
    \begin{tasks}(2)
        \task \( x = \pm \dfrac{2\pi}{3} + k2\pi \)
        \task \( x = \pm \dfrac{\pi}{6} + k\pi \)
        \task \( x = \pm \dfrac{\pi}{3} + k2\pi \)
        \task \( x = \pm \dfrac{\pi}{6} + k2\pi \)
    \end{tasks}
}

% --- CÂU 5 ---
\questionLinesManual{0}{5}{Phương trình \( \cos x = 0 \) có nghiệm là:}{
    \begin{tasks}(2)
        \task \( x = \dfrac{\pi}{2} + k\pi \ (k \in \mathbb{Z}) \)
        \task \( x = k2\pi \ (k \in \mathbb{Z}) \)
        \task \( x = \dfrac{\pi}{2} + k2\pi \ (k \in \mathbb{Z}) \)
        \task \( x = k\pi \ (k \in \mathbb{Z}) \)
    \end{tasks}
}

% --- CÂU 6 ---
\questionLinesManual{0}{6}{Nghiệm của phương trình \( \sin \dfrac{x}{2} = 1 \) là}{
    \begin{tasks}(2)
        \task \( x = \pi + k4\pi, \, k \in \mathbb{Z} \)
        \task \( x = k2\pi, \, k \in \mathbb{Z} \)
        \task \( x = \pi + k2\pi, \, k \in \mathbb{Z} \)
        \task \( x = \dfrac{\pi}{2} + k2\pi, \, k \in \mathbb{Z} \)
    \end{tasks}
}

% --- CÂU 7 ---
\questionLinesManual{0}{7}{Phương trình \( \cos x = \cos \dfrac{\pi}{3} \) có tất cả các nghiệm là:}{
    \begin{tasks}(2)
        \task \( x = \dfrac{2\pi}{3} + k2\pi \, (k \in \mathbb{Z}) \)
        \task \( x = \pm \dfrac{\pi}{3} + k\pi \, (k \in \mathbb{Z}) \)
        \task \( x = \pm \dfrac{\pi}{3} + k2\pi \, (k \in \mathbb{Z}) \)
        \task \( x = \dfrac{\pi}{3} + k2\pi \, (k \in \mathbb{Z}) \)
    \end{tasks}
}

% --- CÂU 8 ---
\questionLinesManual{0}{8}{Tìm tất cả các nghiệm của phương trình \( \cos \dfrac{x}{3} = 0 \).}{
    \begin{tasks}(2)
        \task \( x = k\pi, \, k \in \mathbb{Z} \)
        \task \( x = \dfrac{\pi}{2} + k\pi, \, k \in \mathbb{Z} \)
        \task \( x = \dfrac{3\pi}{2} + k6\pi, \, k \in \mathbb{Z} \)
        \task \( x = \dfrac{3\pi}{2} + k3\pi, \, k \in \mathbb{Z} \)
    \end{tasks}
}

% --- CÂU 9 ---
\questionLinesManual{0}{9}{Phương trình lượng giác: \( \sqrt{3} \cdot \tan x + 3 = 0 \) có nghiệm là:}{
    \begin{tasks}(4)
        \task \( x = \dfrac{\pi}{3} + k\pi \)
        \task \( x = -\dfrac{\pi}{3} + k2\pi \)
        \task \( x = \dfrac{\pi}{6} + k\pi \)
        \task \( x = -\dfrac{\pi}{3} + k\pi \)
    \end{tasks}
}

% --- CÂU 10---
\questionLinesManual{0}{10}{Phương trình \( 2\cot x - \sqrt{3} = 0 \) có nghiệm là}{
    \begin{tasks}(2)
        \task 
        \(\begin{cases} 
        x = \dfrac{\pi}{6} + k2\pi \\ 
        x = -\dfrac{\pi}{6} + k2\pi 
        \end{cases} \quad (k \in \mathbb{Z})\)
        \task 
        \(x = \dfrac{\pi}{3} + k2\pi \quad (k \in \mathbb{Z})\)
        \task 
        \(x = \operatorname{arccot}\dfrac{\sqrt{3}}{2} + k\pi \quad (k \in \mathbb{Z})\)
        \task 
        \(x = \dfrac{\pi}{6} + k\pi \quad (k \in \mathbb{Z})\)
    \end{tasks}
}


\end{document}